\section{Giả thuyết và thiết lập}
\label{sec:hypothesis}

\subsection{Claim có thể bác bỏ (RQ1)}

Chúng tôi phát biểu claim chính để sau đo lường có thể ủng hộ hoặc bác bỏ:

\begin{examplebox}{H1: claim RQ1 có thể bác bỏ}
      Cùng một cách huấn luyện và đánh giá: nếu chuỗi có \textbf{mùa rõ} hoặc \textbf{xu hướng rõ}
      (chỉ số \(F_S\), \(F_T\) từ STL cao \cite{cleveland1990stl,hyndman2021fpp3}), thì
      \[
            \mathrm{MSE}(\mathrm{DLinear}) < \mathrm{MSE}(\mathrm{NLinear})
      \]
      (tức DLinear dự báo tốt hơn NLinear).

      Ngược lại, nếu chuỗi \textbf{hay đổi mức} hoặc \textbf{không dừng mạnh}
      (ví dụ Exchange-Rate, đoạn VIC biến động cao), thì NLinear tốt hơn, hoặc hai model gần nhau
      (\(\lvert\Delta\rvert\) nhỏ).
\end{examplebox}

Định nghĩa vận hành (WEEK02 đến WEEK03):

\begin{align*}
      F_T & = \max\!\bigl(0,\, 1 - \mathrm{Var}(R)/\mathrm{Var}(T+R)\bigr), \\
      F_S & = \max\!\bigl(0,\, 1 - \mathrm{Var}(R)/\mathrm{Var}(S+R)\bigr),
\end{align*}
với \(T,S,R\) lần lượt là xu hướng, mùa và phần dư từ STL. Các proxy không dừng tùy chọn
(\(p\)-value ADF và/hoặc tỷ số mean/std giữa train và test) sẽ được ghi lại khi đủ thời gian làm.
Khoảng cách có dấu là
\(\Delta = \mathrm{MSE}(\mathrm{DLinear}) - \mathrm{MSE}(\mathrm{NLinear})\)
(âm nghĩa DLinear thắng).

\textbf{Quy tắc bác bỏ.}
Nếu các dataset \(F_S\)/\(F_T\) cao \emph{không} cho \(\Delta\) âm một cách có hệ thống,
hoặc các chuỗi dịch mạnh \emph{không} ưu tiên NLinear / khoảng hẹp, ta đánh dấu H1
bị bác bỏ hoặc chưa kết luận, và vẫn công bố pattern đo được.

\subsection{Kỳ vọng phụ (RQ2)}

Ngoài H1, nhóm cũng xem \(L\) (độ dài lookback) ảnh hưởng thế nào tới hai model.
Kỳ vọng: trên chuỗi \textbf{mùa rõ} (\(F_S\) cao), khi tăng \(L\) thì MSE của DLinear và NLinear
lệch nhau rõ hơn, vì DLinear có thể dùng cửa sổ dài để bắt pattern mùa.
Trên chuỗi giá cổ phiếu gần như đi ngẫu nhiên, hai đường MSE theo \(L\) có thể sát nhau hơn.

RQ2 chỉ mang tính khám phá: nhóm sẽ vẽ và mô tả từng dataset, không đặt một claim đúng/sai
cứng như H1.


\subsection{Phác thảo protocol (WEEK02 đến WEEK04)}

\begin{description}[style=nextline,leftmargin=2.2em]
      \item[Mô hình.]
            Cài đặt PyTorch từ đầu cho Linear, NLinear và DLinear;
            channel-independent, trọng số dùng chung giữa các kênh
            \cite{zeng2023ltsf}.
      \item[Metric.]
            MSE và MAE. Với VIC, có thể đánh giá thêm trên returns nếu protocol WEEK03 quy định.
      \item[Lookback / horizon.]
            Mặc định LTSF: \(L,T \in \{96,192,336,720\}\) (có thể cắt bớt và ghi rõ nếu thiếu tài nguyên).
            VIC: \(L \in \{5,30,120,480\}\), độ dài dự báo \(5\).
      \item[WEEK02.]
            Smoke-test mô hình; tính profile STL (ra CSV đặc trưng dữ liệu cho mọi dataset, gồm VIC).
      \item[WEEK03.]
            Bảng train/eval đa dataset theo protocol viết sẵn (split, scaling, seed).
      \item[WEEK04.]
            Ablation kích thước kernel MA của DLinear; đường cong lookback (RQ2);
            ghép kết quả với profile STL cho scatter/heatmap RQ1;
            viết guideline một đoạn và đánh dấu H1 được ủng hộ / bác bỏ / chưa kết luận.
\end{description}

\noindent\textit{Ghi chú chỉnh sửa sau.}
Các mục Methods chi tiết, Results, Guideline và Limitations sẽ được bổ sung vào cùng tài liệu
này khi có thực nghiệm; WEEK02 chủ đích dừng ở khung research-gap ở trên.
